\chapter{Fazit}
\label{ch:fazit}

\section{Erfüllung der Anforderungen}
\label{sec:fazit-anforderungen}

Ausgangspunkt dieser Masterarbeit waren die im Abschnitt~\ref{sec:lps-limitationen} identifizierten Limitationen von \gls{lps}, aus denen in Abschnitt~\ref{sec:alps-anforderungen} die Anforderungen an \gls{alps} abgeleitet wurden. Im Folgenden werden die Anforderungen zusammen mit den Mechanismen, die zur Erfüllung der Anforderung eingesetzt wurden, zusammengefasst.

\begin{itemize}
    \item \textbf{Asynchrone Entkopplung beibehalten}: \gls{alps} übernimmt die gleiche \textit{Proxy-Broker-Proxy-Architektur} von \gls{lps}. Die synchrone Kopplung zwischen \gls{srv}~1 und \gls{srv}~2 wird weiterhin durch die Kette \textit{\gls{alpone} $\rightarrow$ Broker $\rightarrow$ \gls{alptwo}} ersetzt, sodass \gls{alpone} den Eingang einer Alarmmeldung bestätigen kann, sobald sie im Broker gepuffert ist. \gls{alpdedup} wird erst hinter dem Broker in der Verarbeitungskette positioniert und verlängert somit nicht die synchrone Wartezeit der vorgelagerten Systeme. Die Anforderung ist damit architektonisch erfüllt.

    \item \textbf{Transparenz zum Legacy-System beibehalten}: Analog zu \gls{lps} wird \gls{alps} zwischen den bestehenden Services integriert, ohne Modifikationen am Quellcode des Legacy-Systems vorzunehmen. Die Legacy-Schnittstel\-len bleiben nach außen unverändert. Vorausgesetzt wird lediglich, dass Alarmmeldungen die Felder \texttt{device\_id} und \texttt{triggertime} besitzen und wiederholte Alarmmeldungen die gleiche \texttt{triggertime} wie die ursprüng\-li\-che Alarmmeldung tragen. Die Anforderung ist unter dieser Voraussetzung erfüllt.

    \item \textbf{Nachrichtenverlust verbieten}: Alarmmeldungen werden erst dann aus der Queue entfernt, wenn der nachgelagerte Service hinter \gls{alptwo} eine erfolgreiche Verarbeitung bestätigt hat. Im Fehlerfall wird die Alarmmeldung mit \texttt{requeue = true} negativ bestätigt und an ihre originale Position in der Queue zurückgestellt. In allen vier Phasen der experimentellen Evaluation gilt $N_{\text{loss}} = 0$, auch unter den provozierten Netzwerkfehlern in \textit{Phase~4}. Die Anforderung ist somit auf dem Pfad \textit{Broker $\rightarrow$ \gls{alptwo} $\rightarrow$ \gls{srvt}} erfüllt, unter der Voraussetzung, dass der Broker akzeptierte Alarmmeldungen dauerhaft speichert.

    \item \textbf{Exactly-Once-Delivery garantieren}: Die fehlende \textit{Idempotenz} des Legacy-Systems wird durch \gls{alpdedup} ergänzt, das bereits verarbeitete Alarmmeldungen er\-kennt und nicht erneut weiterleitet. Die Fehlerkonstellation, unter der \gls{lps} die \textit{Exactly-Once-Delivery} verletzt hat, wurde in \textit{Phase~1} gezielt reproduziert und führte in keinem der 59 Durchläufe zu einer Mehrfachverarbeitung. Auch in \textit{Phase~3} und \textit{Phase~4} gilt $N_{\text{dup}} = 0$. Da sich die \textit{Exactly-Once-Delivery} in verteilten Systemen nicht uneingeschränkt garantieren lässt~\citep{akkoyunlu1975}, gilt diese Aussage unter idealisierten Bedingungen: Zwischen \gls{alpdedup} und \gls{srvt} treten keine Kommunikationsfehler auf, und \gls{alpdedup}-Instanzen fallen zwischen der Weiterleitung und dem Speichern einer Alarmmeldung nicht aus. Die Anforderung ist unter diesen Bedingungen erfüllt.

    \item \textbf{In-Order-Delivery pro Alarmanlage garantieren}: Die Kombination aus \texttt{x-consistent-hash}, \texttt{x-single-active-\allowbreak consumer} und \texttt{prefetch\_\allowbreak count = 1} stellt sicher, dass Alarmmeldungen einer Alarmanlage immer in derselben Queue landen und sequentiell von genau einem Consumer verarbeitet werden. Durch das Zurückstellen fehlgeschlagener Alarmmeldungen an ihre originale Position bleibt die Reihenfolge auch im Fehlerfall erhalten. Die Fehlerkonstellation, unter der \gls{lps} die \textit{In-Order-Delivery} verletzt hat, führte in \textit{Phase~2} zu keiner Verletzung, und auch unter realitätsnaher Last und Netzwerkfehlern in \textit{Phase~3} und \textit{Phase~4} gilt $N_{\text{ooo}} = 0$. Gleichzeitig können Alarmmeldungen verschiedener Alarmanlagen weiterhin parallel verarbeitet werden. Da Netzwerkfehler oder hohe Latenzen vor dem Broker die Reihenfolge weiterhin verändern können, ist die Anforderung auf dem Pfad \textit{Broker $\rightarrow$ \gls{alptwo} $\rightarrow$ \gls{srvt}} erfüllt.

    \item \textbf{Vendor-Neutralität}: Die verwendeten Mechanismen sind zwar Erweiterungen von \textit{RabbitMQ}, lassen sich funktional jedoch auch auf andere gängige Message Broker wie \textit{Apache ActiveMQ Artemis} oder \textit{LavinMQ} übertragen. Die Anforderung ist damit konzeptionell erfüllt. Die experimentelle Evaluation wurde jedoch ausschließlich mit \textit{RabbitMQ} durchgeführt.
\end{itemize}
